% !TEX program = xelatex
% 在 VSCode 中请使用 XeLaTeX 编译（推荐 TeX Live 或 MiKTeX 完整版，并已安装 ctex 宏包）
\documentclass[12pt,a4paper]{article}
\usepackage[UTF8,scheme=plain]{ctex}
\usepackage{geometry}
\geometry{left=2.5cm,right=2.5cm,top=2.5cm,bottom=2.5cm}
\usepackage{array}
\usepackage{booktabs}
\usepackage{multirow}
\usepackage{xcolor}
\usepackage{makecell}
% 将带圈数字①-⑳声明为CJK字符，使其走中文字体，避免缺字
\xeCJKDeclareCharClass{CJK}{"2460 -> "2473}

% 表格列定义：让内容居中或左对齐更贴近原表
\newcolumntype{C}[1]{>{\centering\arraybackslash}m{#1}}
\newcolumntype{L}[1]{>{\raggedright\arraybackslash}m{#1}}

\linespread{1.3}

\begin{document}

% ============ 标题 ============
\begin{center}
    {\LARGE\bfseries 2026 全国大学生数学建模竞赛}\\[1.2em]
    {\LARGE\bfseries AI 工具使用详情说明}
\end{center}

\vspace{1em}

% ============ 一 ============
\section*{一、所用 AI 工具名称、版本或型号}

\noindent{\kaishu\color{brown}【填写提示】列出竞赛过程中使用的所有 AI 工具，包括大语言模型（如 ChatGPT、豆包、DeepSeek 等）、代码辅助工具、AI 智能体等。版本号要准确填写，不要只写“最新版”。主要用途用一句话概括即可。}

\vspace{0.5em}

\begin{center}
\begin{tabular}{|C{1.2cm}|C{3cm}|C{2.5cm}|L{6.5cm}|}
\hline
\textbf{序号} & \textbf{AI 工具名称} & \textbf{版本/型号} & \multicolumn{1}{C{6.5cm}|}{\textbf{主要用途}} \\
\hline
示例 & ChatGPT & GPT-4o & 模型思路参考、代码调试、文字润色 \\
\hline
1 & & & \\
\hline
2 & & & \\
\hline
3 & & & \\
\hline
4 & & & \\
\hline
\end{tabular}
\end{center}

% ============ 二 ============
\section*{二、具体使用目的和环节}

\noindent{\kaishu\color{brown}【填写提示】按论文写作的 7 个核心环节，逐一说明 AI 工具的具体使用目的。使用目的要具体，不要只写“辅助建模”，应写清 AI 在该环节做了什么（如“提供多种建模思路供本队对比选择”、“辅助编写数据预处理代码”等）。未使用 AI 的环节在“使用目的”栏填写“未使用”。}

\vspace{0.5em}

\begin{center}
\begin{tabular}{|C{1cm}|C{3.2cm}|C{1.5cm}|L{5cm}|C{2.5cm}|}
\hline
\textbf{序号} & \textbf{论文写作环节} & \textbf{是否使用} & \multicolumn{1}{C{5cm}|}{\textbf{使用目的}} & \textbf{使用工具} \\
\hline
1 & 赛题理解与问题分析 & 是/否 & （如：辅助拆解题目要求，识别关键约束条件） & \\
\hline
2 & 模型假设与符号定义 & 是/否 & （如：参考常见假设类型，本队筛选后确定） & \\
\hline
3 & 模型建立与算法设计 & 是/否 & （如：提供算法选型建议，核心建模由本队主导） & \\
\hline
4 & 模型求解与编程实现 & 是/否 & （如：辅助编写求解代码，调试报错，本队验证） & \\
\hline
5 & 结果分析与模型检验 & 是/否 & （如：辅助解读结果，提供检验方法参考） & \\
\hline
6 & 论文撰写与文字润色 & 是/否 & （如：优化文字表达，调整语句通顺度） & \\
\hline
7 & 其他辅助环节（文献/数据/图表等） & 是/否 & （如：辅助整理文献格式，数据清洗） & \\
\hline
\end{tabular}
\end{center}

% ============ 三 ============
\section*{三、主要提示方式与使用过程说明}

\subsection*{（一）主要提示方式}

\noindent{\kaishu\color{brown}【填写提示】说明本队使用 AI 工具的主要交互方式，在“是否使用”栏填写“是”或“否”，并在“简要说明”栏简述使用场景（如“用于代码编写和调试”）。可多选。}

\vspace{0.5em}

\begin{center}
\begin{tabular}{|C{1.2cm}|C{4cm}|C{1.8cm}|L{6cm}|}
\hline
\textbf{序号} & \textbf{提示方式} & \textbf{是否使用} & \multicolumn{1}{C{6cm}|}{\textbf{简要说明}} \\
\hline
1 & 网页对话框交互 & 是/否 & （如：用于问题咨询、思路讨论、文字润色） \\
\hline
2 & 代码编辑器内嵌 AI & 是/否 & （如：用于代码补全、调试、重构） \\
\hline
3 & 上传文件或数据对话 & 是/否 & （如：上传赛题数据或论文草稿进行分析） \\
\hline
4 & AI 智能体工作流（多步自动执行） & 是/否 & （如：使用自动化 Agent 完成数据处理流程） \\
\hline
5 & 其他方式 & 是/否 & （请说明） \\
\hline
\end{tabular}
\end{center}

\subsection*{（二）典型交互示例}

\noindent{\kaishu\color{brown}【填写提示】选取竞赛过程中 2 至 3 次与 AI 的关键交互进行记录。选择标准：①对建模思路或结果有重要影响的交互；②能体现本队如何使用和核验 AI 输出的交互。提示词应完整记录，AI 回复核心内容，本队处理方式要写清楚是直接采纳、修改后采纳还是未采纳，以及具体修改了什么。}

\vspace{0.5em}

% ---- 示例 1 ----
\noindent\textbf{示例 1：}

\begin{center}
\begin{tabular}{|C{4cm}|L{9cm}|}
\hline
\multicolumn{1}{|C{4cm}|}{\textbf{项目}} & \multicolumn{1}{C{9cm}|}{\textbf{填写内容}} \\
\hline
对应环节 & （如：模型建立与算法设计） \\
\hline
使用工具及版本 & （如：ChatGPT GPT-4o） \\
\hline
交互时间 & （如：2026-09-XX XX:XX） \\
\hline
本队提示词 & （完整记录向 AI 提出的问题或指令，可附截图。注意：不要粘贴 API Key、账号密码等敏感信息） \\
\hline
AI 回复内容 & （记录 AI 的回复核心内容，关键部分可原文引用） \\
\hline
本队处理方式 & 直接采纳 / 修改后采纳 / 未采纳（说明：修改了哪些内容？为什么修改？未采纳的原因是什么？） \\
\hline
人工核验方式 & （说明如何确认 AI 回复的正确性：如独立推导验证、代码独立复现、数据交叉核对、文献查证、队员讨论确认等） \\
\hline
\end{tabular}
\end{center}

\vspace{1em}

% ---- 示例 2 ----
\noindent\textbf{示例 2：}

\begin{center}
\begin{tabular}{|C{4cm}|L{9cm}|}
\hline
\multicolumn{1}{|C{4cm}|}{\textbf{项目}} & \multicolumn{1}{C{9cm}|}{\textbf{填写内容}} \\
\hline
对应环节 & （如：模型建立与算法设计） \\
\hline
使用工具及版本 & （如：ChatGPT GPT-4o） \\
\hline
交互时间 & （如：2026-09-XX XX:XX） \\
\hline
本队提示词 & （完整记录向 AI 提出的问题或指令，可附截图。注意：不要粘贴 API Key、账号密码等敏感信息） \\
\hline
AI 回复内容 & （记录 AI 的回复核心内容，关键部分可原文引用） \\
\hline
本队处理方式 & 直接采纳 / 修改后采纳 / 未采纳（说明：修改了哪些内容？为什么修改？未采纳的原因是什么？） \\
\hline
人工核验方式 & （说明如何确认 AI 回复的正确性：如独立推导验证、代码独立复现、数据交叉核对、文献查证、队员讨论确认等） \\
\hline
\end{tabular}
\end{center}

% ============ 四 ============
\section*{四、对 AI 输出的采纳、人工修改和核验的主要情况}

\noindent{\kaishu\color{brown}【填写提示】对 AI 输出内容（语言润色除外）的采纳、人工修改和核验情况进行总体说明。注意：①语言润色类内容无需逐项说明；②核心建模与分析须由本队主导，AI 仅作辅助；③“采纳与修改情况”栏要写清楚是直接采纳、部分修改后采纳还是未采纳，并简述修改内容；④“人工核验方式”栏要写清楚本队如何验证 AI 输出的正确性，不能只写“已核验”。}

\vspace{0.5em}

\begin{center}
\begin{tabular}{|C{1cm}|C{3.5cm}|L{5cm}|L{4.5cm}|}
\hline
\textbf{序号} & \textbf{AI 输出内容类别} & \multicolumn{1}{C{5cm}|}{\textbf{采纳与修改情况}} & \multicolumn{1}{C{4.5cm}|}{\textbf{人工核验方式}} \\
\hline
1 & 建模思路与方法建议 & （如：AI 提供 3 种思路，本队选择第 2 种并改进） & （如：本队独立推导验证，对比文献确认可行性） \\
\hline
2 & 公式推导与理论参考 & & \\
\hline
3 & 代码编写与调试 & & \\
\hline
4 & 结果分析与模型评价 & & \\
\hline
5 & 论文核心论述（摘要/结论等） & & \\
\hline
6 & 其他内容 & & \\
\hline
\end{tabular}
\end{center}

% ============ 核心环节人工主导确认 ============
\section*{核心环节人工主导确认：}

\noindent{\kaishu\color{brown}【填写提示】以下为竞赛核心环节，须由本队独立或主导完成。请逐项确认，在“是否本队主导”栏填写“是”或“否”，并在“本队贡献说明”栏简述本队在该环节的具体工作。}

\vspace{0.5em}

\begin{center}
\begin{tabular}{|C{1.2cm}|C{4.5cm}|C{2.5cm}|L{6cm}|}
\hline
\textbf{序号} & \textbf{核心环节} & \textbf{是否本队主导} & \multicolumn{1}{C{6cm}|}{\textbf{本队贡献说明}} \\
\hline
1 & 模型结构与创新点 & 是/否 & \\
\hline
2 & 公式推导与求解步骤 & 是/否 & \\
\hline
3 & 程序逻辑与参数设置 & 是/否 & \\
\hline
4 & 结果分析与论文核心论述 & 是/否 & \\
\hline
5 & 论文撰写与图表制作等 & 是/否 & \\
\hline
\end{tabular}
\end{center}

\vspace{1.5em}

\noindent\textbf{本队确认：}以上 AI 工具使用情况真实完整，无隐瞒、无虚假陈述。核心建模与分析由本队主导完成，所有 AI 输出内容（语言润色除外）均经过人工审查与核实。如有不实，本队愿承担相应责任。

\end{document}
